\subchapter{Experiment basic crytpographic tools}{Objective: Learn
	how to generate and use basic cryptographic elements: keys, certificates and
	revocation lists}

To demonstrate certificate manipulation, we will create our own Certificate Authority (CA)
and issue two certificates: one for the target device and one for the host computer.
We will demonstrate how these certificates can be used to authenticate their
owners and enable encrypted communication.
We will also look at what happens when a certificate is revocated.

Finally, we will measure asymmetric cryptography performances and compare that
with symmetric cryptography.

\section{Setting up OpenSSL configuration}

The very first step is to create an OpenSSL configuration file.
We could start from an example template configuration available in OpenSSL
source code (<https://github.com/openssl/openssl/blob/master/apps/openssl.cnf>),
but as this is a bit tedious, an example one is provided in the lab folder:
\code{$HOME/__SESSION_NAME__-labs/fundamental-concepts/ca}

We also need to initialize the certificates database (\texttt{index.txt}),
the serial number of the next certificate (\texttt{serial}) and the serial
number of the next CRL (\texttt{crlnumber}).

\begin{bashinput}
$ cd $HOME/__SESSION_NAME__-labs/fundamental-concepts/ca
$ touch index.txt
$ echo 01 > serial
$ echo 1000 > crlnumber
\end{bashinput}

\section{Creating CA key and certificate}

The next step is to create a key pair and certificate for our certificate
authority.
This will allow to issue signed certificates to our host computer and target.

Let's begin by creating an 4096 bits RSA key pair for our CA.
\begin{bashinput}
$ mkdir private
$ openssl genrsa -out private/cakey.pem 4096
\end{bashinput}

All data about the created key can then be shown using:

\begin{bashinput}
$ openssl rsa -text -noout -in private/cakey.pem
Private-Key: (4096 bit, 2 primes)
modulus:
    00:b5:86:2a:ff:6c:6d:4a:33:d7:56:e5:76:5f:be:
    ...
\end{bashinput}

We then will create a certificate corresponding to this key.
This is our CA root certificate, so it is going to be self-signed.
We did not create any configuration file for our PKI, so openssl will
interactively ask for various details of our certificate.

\begin{bashinput}
$ openssl req -new -x509 -days 3650 -config openssl.cnf -key private/cakey.pem -out cacert.pem
You are about to be asked to enter information that will be incorporated
into your certificate request.
What you are about to enter is what is called a Distinguished Name or a DN.
There are quite a few fields but you can leave some blank
For some fields there will be a default value,
If you enter '.', the field will be left blank.
-----
Country Name (2 letter code) [FR]:
State or Province Name (full name) []:
Locality Name (eg, city) [Oullins]:
Organization Name (eg, company) [Bootlin security training]:
Organizational Unit Name (eg, section) []:
Common Name (e.g. server FQDN or YOUR name) []:Training CA
Email Address []:
\end{bashinput}

The generated certificate can then be shown:

\begin{bashinput}
$ openssl x509 -noout -text -in cacert.pem
Certificate:
    Data:
        Version: 3 (0x2)
        Serial Number:
            3f:cd:3a:b8:0f:3b:92:cc:89:ce:f0:ea:a4:ed:5c:90:d7:23:33:ee
        Signature Algorithm: sha256WithRSAEncryption
        Issuer: C=FR, L=Oullins, O=Bootlin security training, CN=Training CA
        Validity
            Not Before: Jan 30 09:34:48 2026 GMT
            Not After : Jan 28 09:34:48 2036 GMT
        Subject: C=FR, L=Oullins, O=Bootlin security training, CN=Training CA
        Subject Public Key Info:
            Public Key Algorithm: rsaEncryption
                Public-Key: (4096 bit)
                Modulus:
                    00:b5:86:2a:ff:6c:6d:4a:33:d7:56:e5:76:5f:be:
                    ...
                Exponent: 65537 (0x10001)
        X509v3 extensions:
            X509v3 Subject Key Identifier:
                5C:33:AA:15:C9:71:3E:FB:DF:0F:A2:9E:65:48:CF:C5:69:C6:14:ED
            X509v3 Authority Key Identifier:
                5C:33:AA:15:C9:71:3E:FB:DF:0F:A2:9E:65:48:CF:C5:69:C6:14:ED
            X509v3 Basic Constraints: critical
                CA:TRUE
    Signature Algorithm: sha256WithRSAEncryption
    Signature Value:
        ae:08:55:8f:83:d7:9e:a2:5d:0a:68:26:73:c2:66:8e:87:5a:
        ...
\end{bashinput}

We now have a CA key and a CA root certificate.
They can be used to generate and verify certificates for other identities in our
PKI.

\section{Creating host computer key and certificate}

Now we will create a key for our host computer and create a certificate signing
request for our CA.

\begin{bashinput}
$ mkdir ../host
$ cd ../host
$ openssl genrsa -out hostkey.pem 4096
$ openssl req -new -key hostkey.pem -out host.csr
You are about to be asked to enter information that will be incorporated
into your certificate request.
What you are about to enter is what is called a Distinguished Name or a DN.
There are quite a few fields but you can leave some blank
For some fields there will be a default value,
If you enter '.', the field will be left blank.
-----
Country Name (2 letter code) [AU]:FR
State or Province Name (full name) [Some-State]:.
Locality Name (eg, city) []:Oullins
Organization Name (eg, company) [Internet Widgits Pty Ltd]:Bootlin security training
Organizational Unit Name (eg, section) []:
Common Name (e.g. server FQDN or YOUR name) []:host computer
Email Address []:

Please enter the following 'extra' attributes
to be sent with your certificate request
A challenge password []:
An optional company name []:
\end{bashinput}

The generated CSR needs to be sent to the CA, in order to obtain the
corresponding certificate.
Here, we have the CA key locally, so we can use it to generate the certificate.

\begin{bashinput}
$ cd ../ca
$ mkdir newcerts
$ openssl ca -config openssl.cnf -notext -in ../host/host.csr -out ../host/host.crt
Using configuration from openssl.cnf
Check that the request matches the signature
Signature ok
Certificate Details:
        Serial Number: 1 (0x1)
        Validity
            Not Before: Feb  5 12:32:01 2026 GMT
            Not After : Feb  5 12:32:01 2027 GMT
        Subject:
            countryName               = FR
            organizationName          = Bootlin security training
            commonName                = host computer
Certificate is to be certified until Feb  5 12:32:01 2027 GMT (365 days)
Sign the certificate? [y/n]:y


1 out of 1 certificate requests certified, commit? [y/n]y
Write out database with 1 new entries
Database updated
\end{bashinput}

We can then show the certificate content:

\begin{bashinput}
$ openssl x509 -noout -text -in ../host/host.crt
Certificate:
    Data:
        Version: 3 (0x2)
        Serial Number: 1 (0x1)
        Signature Algorithm: sha256WithRSAEncryption
        Issuer: C=FR, L=Oullins, O=Bootlin security training, CN=Training CA
        Validity
            Not Before: Feb  5 12:32:01 2026 GMT
            Not After : Feb  5 12:32:01 2027 GMT
        Subject: C=FR, O=Bootlin security training, CN=host computer
        Subject Public Key Info:
            Public Key Algorithm: rsaEncryption
                Public-Key: (4096 bit)
                Modulus:
                    00:ae:59:b2:18:0e:47:cb:8a:6c:db:59:31:20:49:
                    ...
                Exponent: 65537 (0x10001)
        X509v3 extensions:
            X509v3 Subject Key Identifier:
                EA:AE:F9:73:BD:45:3F:6E:6F:39:00:84:AB:D5:CB:70:DC:0E:02:07
            X509v3 Authority Key Identifier:
                76:4D:2F:67:12:41:4B:4E:B6:DD:A5:77:71:E0:43:9A:DF:E9:F0:E1
    Signature Algorithm: sha256WithRSAEncryption
    Signature Value:
        2a:5a:f7:0c:23:5f:42:0c:d0:d1:a3:34:ee:62:c7:f8:2b:40:
        ...
\end{bashinput}

\section{Creating target key and certificate}

Similarly, we will create a key pair for our target and issue the corresponding
certificate.
The best practice is to generate the key directly on the target device. This
reduces the risk of key leakage and is often the only feasible method when keys
are stored in hardware modules.
The certificate will still have to be generated on the CA computer, so the CSR
will have to be transferred between the target and the CA computer.

You can generate the key pair and CSR by running following commands on the
target:

\begin{bashinput}
$ mkdir /keys
$ cd /keys
$ openssl genrsa -out target.key 4096
$ openssl req -new -key target.key -out target.csr
You are about to be asked to enter information that will be incorporated
into your certificate request.
What you are about to enter is what is called a Distinguished Name or a DN.
There are quite a few fields but you can leave some blank
For some fields there will be a default value,
If you enter '.', the field will be left blank.
-----
Country Name (2 letter code) [AU]:FR
State or Province Name (full name) [Some-State]:.
Locality Name (eg, city) []:Oullins
Organization Name (eg, company) [Internet Widgits Pty Ltd]:Bootlin security training
Organizational Unit Name (eg, section) []:
Common Name (e.g. server FQDN or YOUR name) []:target computer
Email Address []:

Please enter the following 'extra' attributes
to be sent with your certificate request
A challenge password []:
An optional company name []:
\end{bashinput}

On the host computer, transfer the CSR file from the target.

\begin{bashinput}
$ cd ..
$ scp root@${BOARD}:/keys/target.csr .
\end{bashinput}

First, as the CA received a CSR request from the target, let's analyze CSR
content to make sure the data are genuine.

\begin{bashinput}
$ openssl req -noout -text -in target.csr
Certificate Request:
    Data:
        Version: 1 (0x0)
        Subject: C=FR, L=Oullins, O=Bootlin security training, CN=target computer
        Subject Public Key Info:
            Public Key Algorithm: rsaEncryption
                Public-Key: (4096 bit)
                Modulus:
                    00:99:95:07:1e:bd:e2:08:6d:02:0e:b7:c6:0c:05:
                    ...
                Exponent: 65537 (0x10001)
        Attributes:
            (none)
            Requested Extensions:
    Signature Algorithm: sha256WithRSAEncryption
    Signature Value:
        98:48:e2:3d:b4:d8:d7:e9:6d:e4:99:94:35:73:99:a5:62:e8:
        ...
\end{bashinput}

As we are happy with the CSR content, we can generate the corresponding
certificate.

\begin{bashinput}
$ cd ca
$ openssl ca -config openssl.cnf -notext -in ../target.csr -out ../target.crt
Using configuration from openssl.cnf
Check that the request matches the signature
Signature ok
Certificate Details:
        Serial Number: 2 (0x2)
        Validity
            Not Before: Feb  5 12:36:57 2026 GMT
            Not After : Feb  5 12:36:57 2027 GMT
        Subject:
            countryName               = FR
            organizationName          = Bootlin security training
            commonName                = target computer
Certificate is to be certified until Feb  5 12:36:57 2027 GMT (365 days)
Sign the certificate? [y/n]:y


1 out of 1 certificate requests certified, commit? [y/n]y
Write out database with 1 new entries
Database updated
\end{bashinput}

You can then transfer this certificate to your target.
You should also transfer the CA certificate in the same move, as we will need it
later.

\begin{bashinput}
$ scp ../target.crt cacert.pem root@${BOARD}:/keys/
\end{bashinput}

\section{Optional: Send a signed and encrypted message (RAW)}

Let's write a message we want to send encrypted from the host computer to the target.

\begin{bashinput}
$ cd ..
$ echo "Hello target, this is a message from the host computer\!" > host_message
$ openssl pkeyutl -in host_message -out host_message.enc -certin -inkey target.crt -encrypt
\end{bashinput}

Note: In real-world applications, RSA is rarely used to encrypt messages
directly.
Instead, it secures a session key for symmetric encryption.
Also, you might see the encrypted message is as the key: with a key size of 4096
bits, the encrypted message uses 512 bytes.

Now let's generate a signature file for this message. We will use the host
private key to generate the signature, recipients will be able to verify this
signature using the host public key contained in the certificate.

\begin{bashinput}
$ openssl dgst -sha256 -sign host/hostkey.pem -out host_message.sign host_message.enc
\end{bashinput}

You can then transfer both the encrypted file and the signature file to the
target along with the host machine certificate.

\begin{bashinput}
$ scp host_message.enc host_message.sign host/host.crt root@${BOARD}:
\end{bashinput}

\section{Optional: Verify a signature and decrypt file (RAW)}

On the target, you can first verify the message is genuine, using the host
certificate.
Here we are using raw messages, so there is no direct command allowing to
validate the message and the whole key chain in one step.
So first, we will verify the host machine certificate and extract the host
public key from it:

\begin{bashinput}
$ cd
$ openssl verify -verbose -CAfile /keys/cacert.pem host.crt
host.crt: OK
$ openssl x509 -pubkey -noout -in host.crt > host_pubkey.pem
\end{bashinput}

Now we can verify the message itself:

\begin{bashinput}
$ openssl dgst -sha256 -verify host_pubkey.pem -signature host_message.sign host_message.enc
Verified OK
\end{bashinput}

And finally, we can decrypt the message:

\begin{bashinput}
$ openssl pkeyutl -in host_message.enc -inkey /keys/target.key -decrypt
Hello target, this is a message from the host computer!
\end{bashinput}

\section{Send a signed and encrypted message (S/MIME)}

Fine, let's write a reply to this message and send it back to the host computer.
But now let's use the S/MIME container format.

\begin{bashinput}
$ echo "Hi host computer, this is my reply message" > target_message
$ openssl smime -sign -in target_message -signer /keys/target.crt -inkey /keys/target.key -out target_message.sign
$ openssl smime -encrypt -in target_message.sign -out target_message.sign.enc host.crt
\end{bashinput}

On the host, you can again transfer the encrypted file back from the target.

\begin{bashinput}
$ scp root@${BOARD}:target_message.sign.enc .
\end{bashinput}

And you can now decrypt it and verify file signature.

\begin{bashinput}
$ openssl smime -decrypt -in target_message.sign.enc -recip host/host.crt -inkey host/hostkey.pem -out target_message.dec -outform smime
$ openssl smime -verify -in target_message.dec -CAfile ca/cacert.pem
\end{bashinput}

\section{Revoke a certificate}

Let's generate a Certificate Revocation List (CRL).
We will start with an empty one.

\begin{bashinput}
$ cd ca
$ openssl ca -config openssl.cnf -gencrl -out ca.crl
\end{bashinput}

You can then verify the CRL content:

\begin{bashinput}
$ openssl crl -in ca.crl -text -noout
Certificate Revocation List (CRL):
        Version 2 (0x1)
        Signature Algorithm: sha256WithRSAEncryption
        Issuer: C=FR, L=Oullins, O=Bootlin security training, CN=Training CA
        Last Update: Feb  5 12:53:56 2026 GMT
        Next Update: Mar  7 12:53:56 2026 GMT
        CRL extensions:
            X509v3 Authority Key Identifier:
                76:4D:2F:67:12:41:4B:4E:B6:DD:A5:77:71:E0:43:9A:DF:E9:F0:E1
            X509v3 CRL Number:
                4096
No Revoked Certificates.
    Signature Algorithm: sha256WithRSAEncryption
    Signature Value:
        71:b0:ca:fa:ae:a7:1b:96:67:98:0c:93:6d:0a:6a:e4:59:d7:
        ...
\end{bashinput}

Now let's pretend the target key was compromised and we want to revoke target
certificate (serial number 02).

\begin{bashinput}
$ openssl ca -config openssl.cnf -revoke newcerts/02.pem
Using configuration from openssl.cnf
Revoking Certificate 02.
Database updated
\end{bashinput}

So far we only revoked this certificate in the PKI internal database
(index.txt):

\begin{bashinput}
$ cat index.txt
V	270205123201Z		01	unknown	/C=FR/O=Bootlin security training/CN=host computer
R	270205123657Z	260205125936Z	02	unknown	/C=FR/O=Bootlin security training/CN=target computer
\end{bashinput}

The CRL can then be re-generated with the updated database. We can again show
it's content to ensure it is correct.

\begin{bashinput}
$ openssl ca -config openssl.cnf -gencrl -out ca.crl
Using configuration from openssl.cnf
$ openssl crl -in ca.crl -text -noout
Certificate Revocation List (CRL):
        Version 2 (0x1)
        Signature Algorithm: sha256WithRSAEncryption
        Issuer: C=FR, L=Oullins, O=Bootlin security training, CN=Training CA
        Last Update: Feb  5 13:02:41 2026 GMT
        Next Update: Mar  7 13:02:41 2026 GMT
        CRL extensions:
            X509v3 Authority Key Identifier:
                76:4D:2F:67:12:41:4B:4E:B6:DD:A5:77:71:E0:43:9A:DF:E9:F0:E1
            X509v3 CRL Number:
                4097
Revoked Certificates:
    Serial Number: 02
        Revocation Date: Feb  5 12:59:36 2026 GMT
    Signature Algorithm: sha256WithRSAEncryption
    Signature Value:
        1f:e2:49:bb:19:26:f2:c0:3b:d3:30:3c:0b:97:8e:ec:76:99:
        ...
\end{bashinput}

Now let's try again to verify the signature of the message we received from the
target:

\begin{bashinput}
$ cd ..
$ openssl smime -verify -in target_message.dec -CAfile ca/cacert.pem
Hi host computer, this my reply message
Verification successful
\end{bashinput}

The verification succeeded because the certificate itself does not indicate
revocation.
To detect revoked certificates, we must explicitly check the Certificate
Revocation List (CRL).
One way to distribute the CRL is to concatenate it with the CA certificate
itself.
Additionally, we have to ask OpenSSL to use these CRL data.

\begin{bashinput}
$ cat ca/cacert.pem ca/ca.crl > ca/cacert_crl.pem
$ openssl smime -verify -in target_message.dec -CAfile ca/cacert_crl.pem -crl_check
Verification failure
40E764B4347F0000:error:10800075:PKCS7 routines:PKCS7_verify:certificate verify error:../crypto/pkcs7/pk7_smime.c:301:Verify error: certificate revoked
\end{bashinput}

\section{Compare performances with symmetric cryptography}

Let's demonstrate the difference of performances between symmetric and
asymmetric cryptography algorithms.
The idea will be to create a big files and encrypt it both using RSA and AES to
compare the encryption and decryption speed.

A first issue will arise: as we have seen both RSA and AES operate on fixed size
blocks of data.
AES operates on 128-bit blocks, while with RSA the size will depend on the key
length.
The solution is then to split the data in several blocks and encrypt these
blocks separately.
For AES this is a common need, so most tools will allow to encrypt long data,
with various block chaining modes.
For RSA, encrypting huge amounts of data is not common: among several reasons,
we will see that RSA performances clearly prevent from doing so.
As a consequence, we will have to create our own software, allowing to encrypt
huge amounts of data with RSA by splitting it in small blocks.

The \texttt{crypto-benchmark.py} script provided in the \texttt{perfs} folder
does that: it allows to encrypt or decrypt files using either RSA with an ad-hoc
block splitting mechanism or AES with ECB or GCM block chaining modes.
Additionally, it allows to benchmark performances all of these operations.

Run this script and observe the results:

\begin{bashinput}
$ cd perfs
$ ./crypto-benchmark.py benchmark ../host/hostkey.pem
\end{bashinput}

